\chapter{Проектирование решения}

\section{Общая идея}

Изучение существующих подходов привело к следующей комбинации
идей, которая и составляет архитектурное ядро настоящей работы.

\textbf{Основная таблица} (далее~--- main) играет роль источника
истины о множестве живых строк. Её обязательная колонка~---
служебный идентификатор \texttt{\_id BIGINT NOT NULL}, монотонно
выдаваемый при \textsc{insert}. В смешанной конфигурации main
дополнительно хранит \emph{закреплённые} (pinned) поля как обычные
физические колонки. Pinned-поля пользователь объявляет явно при
\texttt{CREATE TABLE} как поля, для которых заранее известна
высокая плотность заполнения; они получают полные преимущества
колоночного хранения (сжатие, zone~maps, векторизованные предикаты,
column~pruning).

\textbf{Вспомогательные таблицы} (далее~--- side) создаются по
одной на каждую пару (поле, тип), не объявленную pinned. Для каждой
такой пары заводится двухколоночная таблица вида:
\begin{verbatim}
CREATE TABLE db._dyn_<main>__<field>__<type> (
  _id BIGINT NOT NULL,
  value <type> NOT NULL
);
\end{verbatim}
В side-таблицу записываются только те строки, для которых исходное
поле было заполнено. Идентификатор \texttt{\_id} связывает строку
side-таблицы со строкой main, в которой это поле было встречено.

\textbf{Операция SELECT} реализуется специализированным
физическим оператором \texttt{scan\_computing\_table}, который
сканирует параллельно main и каждую запрошенную side-таблицу,
строит хеш-таблицу \texttt{\_id}~$\to$~позиция-в-выходе по
результатам main, и подставляет значения из side-таблиц в
соответствующие позиции; недостающие \texttt{\_id} остаются
NULL. Регулярные колонки, физически находящиеся в main,
читаются напрямую из её скана~--- без дополнительной операции
слияния. В терминах логической семантики это эквивалентно
LEFT~JOIN main и всех нужных side-таблиц по \texttt{\_id}, но без
создания цепочки JOIN-узлов в плане: операция гибридного сканирования
выполняется одним оператором.

\textbf{Эволюция схемы} происходит на стороне диспетчера: при
\textsc{insert} встретившаяся впервые пара (поле, тип) лениво
порождает соответствующую side-таблицу. Pinned-поля и их типы
зафиксированы при \texttt{CREATE TABLE}; динамическое продвижение
sparse-полей в main в настоящей работе не реализовано~--- секция
``Переход между режимами'' обсуждает, почему этот механизм отнесён
к направлениям будущей работы.

\section{Формализация модели}

Введём обозначения: $T$~--- main-таблица; $N$~--- число строк в
ней; $\mathcal{F}~=~\{f_1, \ldots, f_k\}$~--- множество
JSON-путей, встреченных хотя бы раз; $n_i$~--- число строк, в
которых поле $f_i$ было заполнено ненулевым значением.

Состояние каждого поля $f_i$ определяется одним из двух статусов:
\begin{itemize}
    \item \textbf{Pinned (закреплённое).} Указывается явно при
    \texttt{CREATE TABLE} через параметр \texttt{pinned\_columns}.
    Такое поле хранится в main-таблице как обычная физическая
    колонка, независимо от статистики заполненности. Pinned-поле
    одно для всех типов: при попытке вставить значение
    несовместимого с объявленным типом возвращается ошибка.
    \item \textbf{Sparse (разреженное).} По умолчанию все поля,
    не объявленные pinned, попадают сюда. Поле хранится только в
    отдельной side-таблице \texttt{\_dyn\_<main>\_\_<field>\_\_
    <type>}. Допускается несколько типов одного логического поля:
    каждая пара (поле, тип) получает свою side-таблицу.
\end{itemize}

В системе сохраняется один логический инвариант:
\begin{quote}\itshape
Если main-таблица содержит строку с идентификатором \texttt{\_id},
то на момент видимости этой строки все side-таблицы, описывающие
sparse-поля схемы, уже содержат соответствующие записи для этого
\texttt{\_id} (либо явное значение, либо его отсутствие).
\end{quote}
Из этого инварианта вытекает простое правило чтения: оператор
\texttt{scan\_computing\_table} использует main как ``якорь'' для
определения живых строк; видимая строка main гарантированно имеет
консистентное представление в side-таблицах. Side-таблицы могут
содержать ``висящие'' записи, не отражённые в main (например, из-за
обрыва транзакции между записью side и записью main),~--- такие
записи никогда не попадут в результат \textsc{select}, поскольку
\texttt{\_id} отсутствует в main и игнорируется при гибридном
сканировании. Подробное обсуждение этого инварианта и его роли
в обеспечении параллельности приведено в разделе
``Параллельность''.

\section{Обоснование формы хранения}

\subsection{Отдельная таблица вместо внутреннего формата}

Центральное архитектурное решение~--- хранить разреженные поля в
полноценных двухколоночных таблицах, а не в специальном внутреннем
формате (как \texttt{shared~data} в ClickHouse). Аргументы в пользу
этого выбора:

\begin{itemize}
    \item \textbf{Переиспользование всей реляционной машинерии.}
    Любая оптимизация, применяемая к основной таблице
    (zone~maps, векторизация, projection~pushdown, индексы,
    сжатие), автоматически применима и к вспомогательной таблице.
    Не требуется писать специализированный код для разреженного
    хранилища.
    \item \textbf{Совместимость с MVCC и WAL.} Операции над
    вспомогательной таблицей проходят по тому же пути, что и над
    основной, автоматически получая транзакционные гарантии и
    журналирование.
    \item \textbf{Простота отладки и диагностики.}
    Вспомогательная таблица~--- это обычная таблица; её можно
    посмотреть отдельным \textsc{select}, проверить содержимое,
    оценить размер, применить EXPLAIN.
    \item \textbf{Возможность независимой оптимизации.}
    Вспомогательная таблица может иметь собственный индекс по
    \texttt{\_id} (например, zone~map по \texttt{\_id} полезен
    при выборочных предикатах по основной таблице с последующим
    JOIN).
\end{itemize}

\subsection{Логический идентификатор вместо физической позиции строки}

Альтернатива~--- использовать физический номер строки
(row~position) как ключ связи. Отбрасывается по причинам
устойчивости к операциям над основной таблицей (физическая позиция
может меняться при сжатии, ребалансировке row~groups, перестройке
индексов; логический \texttt{\_id} остаётся стабильным) и
отсутствия кросс-зависимостей (вспомогательная таблица не должна
знать о физическом расположении строк в основной; это упрощает
параллельную обработку~--- \textsc{insert} в основную и
вспомогательные таблицы могут идти одновременно, как только известен
\texttt{\_id}).

Идентификаторы \texttt{\_id} генерируются монотонно: для каждой
операции \textsc{insert} система запрашивает у основной таблицы
текущее число строк, присваивает вставляемым строкам
последовательный диапазон $[\texttt{start\_id},
\texttt{start\_id}+|C|)$, и этот диапазон далее используется как
ключ связи.

\subsection{Специализированный оператор гибридного сканирования}

Чтение данных с разреженными полями реализовано отдельным физическим
оператором \texttt{scan\_computing\_table}, а не цепочкой LEFT~JOIN
в логическом плане. Семантически оба варианта эквивалентны, но
выделение в отдельный оператор имеет ряд преимуществ:

\begin{itemize}
    \item \textbf{Локальность алгоритма.} Параллельное чтение
    main и всех нужных side-таблиц, построение хеш-таблицы
    \texttt{\_id}~$\to$~позиция и подстановка значений происходят
    внутри одного оператора. Это позволяет не материализовывать
    промежуточные представления отдельных таблиц.
    \item \textbf{Естественная интеграция с column~pruning.}
    Оператор принимает список запрошенных колонок и обращается
    только к тем side-таблицам, поля которых упомянуты в запросе.
    Пары (поле, тип), не упомянутые в запросе, не сканируются
    вовсе.
    \item \textbf{Единый источник истины.} Множество <<живых>>
    \texttt{\_id} определяется одним сканом main; для каждой
    side-таблицы строится отдельный проход с подстановкой
    значений по уже построенной хеш-таблице. Не возникает
    проблем согласованности нескольких независимых
    JOIN-операторов.
    \item \textbf{Совместимость с параллельностью.} Из инварианта
    раздела ``Формализация'' следует, что сканирование main и
    side-таблиц можно вести параллельно, не дожидаясь окончания
    параллельных \textsc{insert}: видна будет только та часть
    данных, которая зафиксирована в main на момент начала
    транзакции (см. раздел ``Параллельность'').
\end{itemize}

Альтернатива в виде LEFT~JOIN на уровне логического плана была
отброшена: она вынуждает оптимизатор и валидатор учитывать
служебные узлы JOIN, что увеличивает поверхность изменений в
зрелых компонентах кодовой базы, не давая дополнительной
функциональности~--- алгоритмически операция остаётся такой же.

\section{Когда стоит делать колонку pinned}

В предложенной модели решение о размещении поля принимается
однократно на этапе \texttt{CREATE TABLE} и далее не меняется.
Аналитик объявляет некоторый список pinned-полей (возможно, пустой);
все остальные поля автоматически становятся sparse. Этот выбор
влияет на три аспекта производительности:

\begin{itemize}
    \item \textbf{Скорость \textsc{select} по полю.} Чтение
    pinned-поля происходит напрямую из колоночного буфера
    main-таблицы и пользуется всей колоночной машинерией
    (zone~maps, векторизация, сжатие). Чтение sparse-поля
    требует прохода по соответствующей side-таблице с
    хеш-подстановкой по \texttt{\_id}~--- это дороже на
    константный множитель, но всё ещё линейно по числу строк.
    \item \textbf{Стоимость хранения редких полей.} Pinned-поле
    выделяет в main-таблице буфер на все $N$ строк независимо от
    плотности заполнения. Для редкого поля это пустая трата
    памяти. Sparse-поле занимает место только под фактически
    заполненные строки.
    \item \textbf{Стоимость эволюции схемы.} Pinned-поля
    фиксированы при \texttt{CREATE TABLE}; добавление нового
    pinned-поля задним числом не предусмотрено. Sparse-поля
    добавляются лениво в момент первой вставки без блокировок и
    миграции.
\end{itemize}

Эмпирическое правило: если поле упоминается в большинстве запросов
и заполнено в более чем 30--50\% строк~--- кандидат на pinned. Поля
с плотностью ниже 10\% разумно оставлять sparse независимо от
частоты использования.

Альтернативный механизм автоматического продвижения sparse-поля в
main по достижении порога заполненности рассматривался на ранней
стадии работы и отнесён к будущим направлениям; обсуждение
см.~раздел ``Переход между режимами''.

\section{Переход между режимами}

Автоматическое продвижение sparse-поля в pinned-колонку
main-таблицы по достижении порога заполненности рассматривалось на
ранней стадии работы. В предложенной архитектуре от него решено
отказаться по совокупности причин.

\textbf{Сложность под нагрузкой.} Продвижение~--- $O(N)$ операция
(добавление физической колонки в main с инициализацией значений по
умолчанию + полное сканирование side-таблицы для миграции
значений). Если продвижение запускается в фоне параллельно с
\textsc{insert}, требуется аккуратная синхронизация: новые
\textsc{insert} могут писать в side, ещё не дочитанную миграцией.
Если выполняется блокирующим образом, продвижение приостанавливает
все \textsc{insert} в данное поле на время операции.

\textbf{Нарушение инварианта во время миграции.} Главное достоинство
текущей схемы~--- стабильное правило ``main определяет видимость, side
дополняет значения''. Промежуточное состояние ``часть строк уже в
main-колонке, часть ещё в side-таблице'' нарушает этот инвариант и
требует двойного источника правды на время миграции, что
существенно усложняет реализацию \texttt{scan\_computing\_table}.

\textbf{Необходимость обратного перехода.} Если продвижение
осуществляется автоматически по статистике, симметрично следует
поддерживать и понижение (удаление физической колонки при падении
плотности). Эта операция $O(N)$ и блокирующая по той же причине,
что и прямой переход.

В предложенной архитектуре выбран альтернативный путь: статус поля
(pinned либо sparse) фиксируется при \texttt{CREATE TABLE} и в
дальнейшем не меняется. Это позволяет сохранить инвариант
``main~--- источник истины'' и упростить и оператор сканирования, и
анализ параллельной семантики (см.~раздел ``Параллельность'').
Автоматическое продвижение рассматривается как направление будущей
работы в случаях, когда статистика заполненности неизвестна на
этапе \texttt{CREATE TABLE} и заранее зафиксировать список pinned
невозможно.

\section{Плюсы и минусы подхода}

\textbf{Преимущества:}
\begin{enumerate}
    \item Колоночная производительность на pinned-полях: они
    физически лежат в main и пользуются всей колоночной
    машинерией (zone~maps, векторизация, сжатие, column~pruning).
    \item Компактное хранение редких полей: side-таблицы содержат
    только фактически заполненные строки; неупомянутые в запросе
    side-таблицы не сканируются вовсе.
    \item Автоматическая эволюция схемы: не требуется ручного
    переопределения структуры при появлении новых sparse-полей~---
    side-таблица для новой пары (поле, тип) создаётся лениво при
    первой вставке.
    \item Отсутствие специальных структур данных: side-таблица~---
    это обычная колоночная таблица; вся транзакционная и
    инфраструктурная машинерия (MVCC, WAL, индексы) применима
    без модификаций.
    \item Стабильная семантика \texttt{\_id}: внешние индексы и
    ссылки на строки устойчивы к внутренним операциям движка
    хранения.
    \item Простой инвариант параллельности (раздел
    ``Параллельность''): если строка видна в main, её
    представление в side-таблицах гарантированно консистентно.
\end{enumerate}

\textbf{Недостатки и компромиссы:}
\begin{enumerate}
    \item Чтение sparse-поля дороже чтения pinned-колонки:
    хеш-подстановка по \texttt{\_id} вносит накладные расходы
    порядка $O(n_i)$ операций хеширования на запрос.
    \item Отсутствие автоматического продвижения: при плохом
    выборе списка pinned-полей часто запрашиваемые
    sparse-поля сохраняют повышенную стоимость чтения. Решение
    отнесено к направлениям будущей работы.
    \item Число side-таблиц может расти значительно: при
    10\,000 уникальных путей в каталоге возникает 10\,000
    физических таблиц; каждая потребляет некоторый объём
    памяти на метаданные.
\end{enumerate}

\section{Псевдокод INSERT}

Для формализации опишем \textsc{insert} в таблицу с вычисляемой
схемой в виде псевдокода (листинг~\ref{lst:insert}). Алгоритм
выполняется одним физическим оператором
\texttt{operator\_insert\_computing} в исполнителе под одной
транзакцией, что гарантирует атомарность по отношению к
параллельным запросам.

\begin{algorithm}[!h]
\caption{Алгоритм \textsc{insert} в таблицу с вычисляемой схемой}
\label{lst:insert}
\begin{algorithmic}
  \Function{Insert}{row-chunk $C$, схема $S$, main-таблица $T$}
    \State \texttt{// pre-flight DDL на стороне диспетчера}
    \ForAll{колонок $c_j$ из $C$}
      \State $S.\mathrm{append}(f_j, t_j)$
      \If{пара $(f_j, t_j)$ ещё не имеет side-таблицы и не
        pinned}
        \State создать side-таблицу
          \texttt{\_dyn\_<T>\_\_<$f_j$>\_\_<$t_j$>}
      \EndIf
    \EndFor
    \State \texttt{// выполнение в исполнителе под одной транзакцией}
    \State $\mathtt{start\_id} \gets$ число строк $T$
    \State $\mathtt{main\_chunk} \gets [\_id]$ + значения
      pinned-колонок из $C$
    \ForAll{колонок $c_j$ из $C$, не являющихся pinned}
      \State построить $\mathtt{side\_chunk}_j$ из пар
        $(\mathtt{start\_id}+i, c_j[i])$ по ненулевым строкам
    \EndFor
    \State \texttt{// первая фаза: side-таблицы}
    \ForAll{$\mathtt{side\_chunk}_j$}
      \State \texttt{storage\_append} в side-таблицу
        $(f_j, t_j)$ с $\mathit{txn\_id}$
    \EndFor
    \State \texttt{// вторая фаза: main-таблица}
    \State \texttt{storage\_append} в $T$ с
      $\mathtt{main\_chunk}$ и $\mathit{txn\_id}$
    \State \texttt{// фиксация}
    \State $\mathit{commit\_id} \gets
      \mathrm{txn\_manager}.\textsc{commit}()$
    \State \texttt{storage\_commit\_append} для каждой
      затронутой таблицы (main и side) с $\mathit{commit\_id}$
  \EndFunction
\end{algorithmic}
\end{algorithm}

\textbf{Принципиальный порядок записи: side первыми, main
последней.} Это поддерживает инвариант ``main определяет
видимость, side хранит данные''. Если транзакция завершилась
успешно, все side-записи с $\mathit{commit\_id}$ становятся видимыми
одновременно с main-записью с тем же $\mathit{commit\_id}$~--- их
снимки атомарны. Если транзакция оборвалась после записи в
side-таблицу, но до записи в main, side-записи остаются помечены
исходным $\mathit{txn\_id}$ (без $\mathit{commit\_id}$) и невидимы
любому читателю. Поскольку \texttt{\_id} такой строки в main также
отсутствует, оператор \texttt{scan\_computing\_table} никогда не
обратится к этой строке, даже если бы видел её.

Алгоритмическая модель \textsc{select} симметрична: один
\texttt{storage\_scan} по main даёт множество живых \texttt{\_id} и
значения pinned-колонок, дальше каждая упомянутая в запросе
sparse-side сканируется один раз с подстановкой значений в
позиции, найденные по хеш-таблице
\texttt{\_id}~$\to$~позиция-в-выходе.

\section{Параллельность}\label{sec:parallel}

Корректность параллельного выполнения операций над таблицей с
вычисляемой схемой обеспечивается не на уровне внешней сериализации
(хеш-маршрутизацией ``одна коллекция $\to$ один исполнитель'',
которая на текущем этапе используется в реализации, но в будущем
будет ослаблена для повышения масштабируемости), а на уровне самого
устройства схемы: \emph{инвариант ``main~--- источник истины''}
делает параллельное чтение и запись безопасными независимо от
порядка планирования транзакций.

\subsection{Базовый инвариант и его следствия}

Напомним инвариант, введённый в разделе ``Формализация модели'':
\begin{quote}\itshape
Если строка с идентификатором \texttt{\_id} видна читателю в main,
то соответствующие записи во всех side-таблицах либо уже видны тому
же читателю, либо явно отсутствуют (что трактуется как NULL).
\end{quote}

Инвариант поддерживается на этапе записи через порядок
\emph{side $\to$ main}: исполнитель сначала проводит
\texttt{storage\_append} во все side-таблицы под текущим
$\mathit{txn\_id}$, затем~--- в main, и только потом получает у
менеджера транзакций общий $\mathit{commit\_id}$ и применяет
\texttt{storage\_commit\_append} ко всем затронутым таблицам с этим
$\mathit{commit\_id}$. MVCC-протокол хранилища гарантирует, что
никакой читатель не видит записи с $\mathit{txn\_id}$ (только
$\mathit{commit\_id}$, выданный после успешного коммита), поэтому
запись становится видимой во всех таблицах одновременно.

При оборванной транзакции (отказ процесса, явный
\texttt{rollback}) ни одна из записей не получает
$\mathit{commit\_id}$~--- они остаются невидимыми. ``Висящие''
записи в side-таблицах безопасно игнорируются последующим
\texttt{scan\_computing\_table}, поскольку без соответствующей
живой строки в main их \texttt{\_id} не попадает в множество ключей
для хеш-подстановки.

\subsection{Параллельный INSERT}

Рассмотрим две параллельные транзакции $\tau_A$ и $\tau_B$, каждая
выполняющая \textsc{insert} нескольких строк в одну и ту же
вычисляемую таблицу.

Каждая транзакция вначале запрашивает у менеджера хранилища текущее
число строк main как $\mathtt{start\_id}$. Эта операция атомарна на
уровне акторного менеджера хранилища: его mailbox обрабатывает
сообщения последовательно, поэтому два параллельных запроса
получают разные значения, отражающие фактический порядок их
обработки. Соответственно $\tau_A$ и $\tau_B$ получают
непересекающиеся диапазоны \texttt{\_id}.

Дальше каждая транзакция действует независимо. Записи $\tau_A$ в
side-таблицы помечены $\mathit{txn\_id}_A$ и невидимы $\tau_B$ (и
наоборот). Когда $\tau_A$ применяет \texttt{storage\_commit\_append}
со своим $\mathit{commit\_id}_A$, все её записи становятся видимыми
последующим транзакциям с большим $\mathit{start\_time}$. Записи
$\tau_B$, ещё не закоммиченные, остаются невидимыми. Никакого
взаимодействия между $\tau_A$ и $\tau_B$ нет: они работают с
непересекающимися \texttt{\_id} и не делят ни одной строки.

\subsection{Параллельный SELECT и INSERT}

Параллельный \textsc{select} в момент незавершённой $\tau_A$
открывает свою транзакцию $\tau_S$ с собственным
$\mathit{start\_time}$. Сканирование main отдаёт только те строки,
$\mathit{commit\_id}$ которых не превышает $\mathit{start\_time}$
читателя~--- стандартный snapshot isolation. Если $\tau_A$ ещё не
закоммитилась к моменту начала $\tau_S$, её записи не видны ни в
main, ни в side. Если $\tau_A$ закоммитилась раньше начала
$\tau_S$, её записи видны и там, и там одновременно~--- инвариант
``main определяет видимость''.

Случай, когда $\tau_A$ закоммитилась в момент исполнения $\tau_S$
(между сканом main и сканом side), требует уточнения. MVCC-снимок
фиксируется при \texttt{begin\_transaction}, не при каждом
обращении к хранилищу,~--- поэтому даже если физически сканы main и
side происходят в разное время, оба видят одно и то же состояние,
консистентное с $\mathit{start\_time}$ читателя. ``Подоспевшие''
изменения $\tau_A$ не попадают в результат $\tau_S$, что и требуется
от snapshot~isolation.

\subsection{Параллельный DELETE}

\textsc{delete} строк с заданным WHERE реализуется отдельным
оператором \texttt{operator\_delete\_computing} и выполняется в
строго обратном порядке: \emph{main первым, side после}.

Сначала исполнитель помечает в main все строки, попадающие под
WHERE, как удалённые (под $\mathit{txn\_id}$); затем~--- те же
строки в каждой side-таблице. Когда транзакция коммитится, помеченные
``удалённые'' записи получают $\mathit{commit\_id}$ и перестают быть
видимы новым читателям. Симметрия с \textsc{insert}: видимость
``удалённого состояния'' возникает одновременно во всех таблицах.

Параллельный \textsc{select} в момент незавершённого DELETE по
snapshot~isolation видит либо все строки целиком (если его
$\mathit{start\_time}$ предшествует $\mathit{commit\_id}$ удаления),
либо ни одной из удаляемых (после коммита). Никаких ``половинных''
состояний~--- то есть видимая в main строка без значений из side
или, наоборот, висящие в side значения без живой строки в main~---
никогда не наблюдается.

Параллельный \textsc{insert} в момент незавершённого DELETE
работает с непересекающимися \texttt{\_id}, поскольку
\texttt{start\_id} для нового \textsc{insert} больше любого
существующего \texttt{\_id}, а удаляемые \texttt{\_id}
заведомо меньше.

\subsection{Параллельный UPDATE}

\textsc{update} реализован как ``re-alloc с полным копированием'':
оператор \texttt{operator\_update\_computing} вычитывает строки,
попадающие под WHERE, применяет SET-выражения к их значениям,
выделяет новый диапазон \texttt{\_id} через
\texttt{storage\_total\_rows}, делает \textsc{insert} в side и main
с новыми \texttt{\_id} и \textsc{delete} старых строк в main и
side. Всё это в одной транзакции.

Параллельный \textsc{select} в момент незавершённого UPDATE по
снимку видит либо старую версию строки (если до коммита), либо
новую (если после). Никогда не видит обе одновременно, поскольку и
вставка ``нового'' и удаление ``старого'' получают один
$\mathit{commit\_id}$.

Параллельные UPDATE/INSERT/DELETE по \emph{разным} \texttt{\_id}
работают полностью независимо: операции над диапазонами с разными
\texttt{\_id} не делят ни одной строки и потому не имеют конфликтов
данных.

Параллельные UPDATE по \emph{одному и тому же} логическому
\texttt{\_id} (с пересекающимися WHERE-условиями) образуют
write-write конфликт. Otterbrix~--- как и большинство классических
аналитических колоночных систем~--- не обнаруживает такие конфликты
автоматически; разрешение в таких случаях соответствует
last-writer-wins по $\mathit{commit\_id}$. Это ограничение
существующей MVCC-реализации, не связанное со схемой вычисляемой
таблицы; устранение требует доработки на уровне всего хранилища.

\subsection{Безопасность по отношению к будущему ослаблению
сериализации}

Текущая реализация хеш-маршрутизирует запросы к одной коллекции на
один акторный исполнитель, что само по себе сериализует все
операции с этой коллекцией. Однако такое ограничение является
артефактом текущего архитектурного решения, не свойством схемы
вычисляемой таблицы. В будущем планируется балансировать запросы по
исполнителям независимо от коллекции (например, по сессии или
по размеру очереди), что увеличит масштабируемость, но потребует
полагаться только на свойства самой схемы для корректности.

Все приведённые выше рассуждения о параллельности \emph{не зависят}
от хеш-маршрутизации: они полагаются исключительно на
MVCC-механизмы хранилища и порядок записи ``side $\to$ main''
(симметрично для удаления ``main $\to$ side''). Таким образом,
предложенная архитектура естественно совместима с будущим
ослаблением сериализации без дополнительных модификаций.

\chapterconclusion

Предложенная модель хранения сочетает сильные стороны колоночного
представления (для часто используемых pinned-полей) и разреженного
EAV (для редких полей), обходя ключевые ограничения каждого из
подходов. Центральные архитектурные решения~--- хранение
sparse-полей в отдельных двухколоночных таблицах, связанных через
стабильный логический идентификатор \texttt{\_id}, гибридное
сканирование одним специализированным оператором, и инвариант
``main~--- источник истины'' для согласованности параллельных
операций~--- позволяют достичь приемлемого баланса между стоимостью
хранения, стоимостью чтения и сложностью реализации, а также
естественной совместимости с будущим ослаблением сериализации на
уровне исполнителей.
